\section{简易程序、公益诉讼程序}
\begin{enumerate}
  \item 适用于\zd{基层法院及派出法庭}审理的\zd{第一审}（\bdy 一审）民事案件。
  \item 不许适用简易程序。
        \begin{enumerate}
          \item 原、被告\zd{送达地址不明确} \yjt 简易程序\zd{不允许公告送达} \yjt 没法送达文书没法审理 \yjt \zd{裁定驳回起诉}。
          \item 重审、再审、代表人诉讼、公益诉讼、第三人撤销之诉。
        \end{enumerate}
  \item 依职权 \hz 协议适用。普通可约成简易 + 法院同意，\zd{简易不准约成普通}（浪费司法资源）。
  \item \zd{传唤、送达方式简单}。
        \begin{enumerate}
          \item \zd{简便方式}送达\zd{所有}文书且\zd{无需同意}。\al{发邮件/传真/短信、打电话、捎口信、农村高音喇叭喊}。
          \item \zd{开庭通知}要确保当事人\zd{收悉}，\zd{不允许缺席判决}。\al{甲起诉乙，一审适用简易，电话通知双方，乙未到庭，法院缺席判决}。
        \end{enumerate}
  \item 程序。\zd{受理、送达、审理、裁判文书}均可以简化。\al{庭审简化、不开庭判决 \cuo}。
        \begin{enumerate}
          \item \zd{先行调解}。婚姻家庭、继承、劳动合同、交通事故、工伤、宅基地、相邻关系、合伙合同、小额诉讼。
          \item 独任制，庭审方式灵活。
          \item 举证期限、答辩期限协商 \hz 法院确定。不得超过 15 天。
          \item 审期 3 + 1，院长批准。
          \item 程序的转换。
                \begin{enumerate}
                  \item 普通转简易。\zd{开庭后}不允许转为简易程序。
                  \item 简易转普通。\zd{用裁定}；审期不重新计算。
                \end{enumerate}
        \end{enumerate}
  \item \zd{小额诉讼的程序} = 简易程序 + 简单金钱给付。
        \begin{enumerate}
          \item 金额标准。\zd{该省上年度就业人员平均工资}为基准来计算。
                \begin{enumerate}
                  \item < 一半 \yjt 依职权。
                  \item 一半到 2 倍 \yjt 协议。
                  \item > 2 倍 \yjt \zd{不准用}。
                \end{enumerate}
          \item \zd{海事案件符合条件，海事法院也可以适用}。
          \item 不允许适用（\zd{小额一审终审，剥夺上诉权利}）。人身关系、财产确权、涉外、诉前鉴定结果有异议（争议较大）、当事人下落不明、\zd{反诉}。
          \item 举证期限协商 \hz 法院确定。一般不得超过 7 天。
          \item \zd{能一次开庭审结就当庭宣判}，不能就多搞几次（\al{应当当庭宣判 \cuo}）；审期 2 + 1，院长批准。
          \item \zd{一审终审}，不能上诉。
          \item 判决简化。认定事实、裁判理由\zd{可以省略}。
          \item 转。遇到最先能转的那个。\al{甲提出了反诉}。
          \item 再审。
                \begin{enumerate}
                  \item 可以向\zd{原审法院}申请再审；合议庭（没有提审）。
                  \item 再审判决能否上诉看\zd{当事人申请再审的理由}。
                        \begin{enumerate}
                          \item 以原判决的事实、法律、程序、枉法裁判为由 \yjt 继续适用小额诉讼一审终审。
                          \item 以不应按照小额诉讼的程序审理为由 \yjt 不再适用一审终审的规定，可以上诉。
                        \end{enumerate}
                \end{enumerate}
        \end{enumerate}
  \item 公益诉讼。仅限组织、机构，个人无法当原告。
        \begin{enumerate}
          \item 起诉条件加一个\zd{初步证据}。
          \item 侵权行为地 \hz 被告住所地中院。
          \item 受理后，\zd{法院应当10 天内}书面告知行政主管部门（通过行政手段效果更好）。
          \item \zd{法院的释明}。
                \begin{enumerate}
                  \item 原告的诉讼请求不足以维护公共利益时，法院告诉原告，来，我告诉你该怎么告。
                  \item 原告坚持不增加、变更，应当以原诉讼请求，作为裁判依据，不能主动增加诉讼请求 \yjt 处分权。
                \end{enumerate}
          \item 其他机关、组织参与诉讼。\al{A 起诉化工厂。BC 开庭前共同原告身份来（\scx{合并审理}），另行起诉，不予受理}。
          \item 可以和解、调解。
                \begin{enumerate}
                  \item \zd{达成和解协议之后，不允许撤诉结案}。
                  \item 必须要公告，期满出具调解书结案 \yjt 唯一结案方式。
                \end{enumerate}
          \item 不允许反诉。
          \item 检察机关提起的公益诉讼。考纲没删除。
                \begin{enumerate}
                  \item 不能成为常态，可能会破坏平等原则。
                  \item 必须满足\zd{补充性}。
                        \begin{enumerate}
                          \item \zd{诉前通告} \yjt 没有法律规定的其他机关组织 \yjt 检察机关。
                          \item 侵害烈士 \yjt 英雄烈士的近亲属 \yjt 公告期满 \yjt 检察机关。
                        \end{enumerate}
                  \item 中院所对应的市检。公益诉讼起诉人 = 原告。
                \end{enumerate}
        \end{enumerate}
\end{enumerate}
